\documentclass[10pt,a4paper]{amsart}
\usepackage[margin=19mm]{geometry}
\usepackage{amsmath,amssymb,mathtools}
\usepackage[scr=rsfs,cal=euler]{mathalfa}
\usepackage{xfrac}
\usepackage[unicode,hidelinks,bookmarks=false]{hyperref}
\newcommand{\ie}{i.e.\ }
\newcommand{\cf}{cf.\ }
\newcommand{\resp}{respectively\ }
\newcommand{\Subsection}{\subsection}
\providecommand{\nobreakdash}{\nobreak\hbox{-}}
\setlength{\emergencystretch}{2em}
\multlinegap0pt
\usepackage{bm}
\usepackage{bbm}
\usepackage{dsfont}
\usepackage{xspace}
\usepackage{cancel}
\usepackage{mathtools}
\usepackage{amsthm}
\makeatletter
\@ifundefined{thm}{\newtheorem{thm}{Theorem}}{}
\@ifundefined{definition}{\newtheorem{definition}[thm]{Definition}}{}
\@ifundefined{lemma}{\newtheorem{lemma}[thm]{Lemma}}{}
\makeatother
\newcommand{\diam}{\operatorname{diam}}
\title[Finite dimensional amenable groups]{Finite dimensional amenable groups}
\author{Anna Erschler, Ivan Mitrofanov}
\date{Source: arXiv:2508.20296v1 (CC BY 4.0)}
\begin{document}
\maketitle

\noindent\emph{Source and licence.} Finite dimensional amenable groups. Version arXiv:2508.20296v1 (CC BY 4.0), distributed under the Creative Commons Attribution 4.0 International licence, as recorded in the arXiv licence metadata for this exact version and shown on the abstract page https://arxiv.org/abs/2508.20296v1. Full rights notice: licenses/arxiv-2508.20296-CC-BY-4.0.txt.\par\smallskip

\noindent\textit{Editorial scope.} Counted unit: the Thm whose source label is \texttt{pr:dimtoisodiam} in the article (source0.tex lines 479--484, source label \texttt{pr:dimtoisodiam}), delivered together with the article's own complete original proof of it (source0.tex lines 486--607, one \texttt{proof} environment of 122 lines). No mathematics is written, repaired or re-derived: statement and proof are the article's own text line for line. Required context retained uncounted: defn:couples (lines 397--406). Every definition, display and label of those lines is retained unchanged. Omitted from this package and not redistributed: 1--396, 407--478, 485--485, 608--1237. Nothing inside a retained excerpt is omitted or altered: every retained body is delivered exactly as the article's lines are written, so no omission receipt of any kind accompanies this package. Citations: the retained excerpts carry no citation command at all, so no bibliography entry of the article is transcribed and no reference is redistributed. Reference adaptation: none was needed and none was made; every reference inside the delivered unit resolves inside it, so no unresolved reference or citation is produced. Printed slips kept literal: no printed word, symbol, hypothesis or number of the counted statement or of its proof was altered. Layout and class: the article's own class is \texttt{amsart} with packages including enumerate, amssymb, amsmath, graphicx, ucs, tikz, mathtools, wrapfig, hyperref, todonotes, inputenc, times, mathabx, amsfonts; the wrapper is the lane's pinned \texttt{amsart} editorial wrapper at 10pt on A4, which restates the theorem-like environments the retained text needs and adds the article's own macro definitions that the retained text needs (\texttt{\textbackslash diam}), with extra wrapper packages (bm, bbm, dsfont, xspace, cancel, mathtools). The delivered numbering is the unit's own. Assets: no retained span contains a figure environment, a graphics-inclusion command, a PDF-page inclusion, a table or a TikZ picture; no image file is copied into this package, the package declares no asset dependency and no image file is redistributed with it. Licence: version arXiv:2508.20296v1 of this article is distributed under the Creative Commons Attribution 4.0 International licence, as recorded in the arXiv licence metadata for this exact version and shown on the abstract page https://arxiv.org/abs/2508.20296v1. A full rights notice is delivered at licenses/arxiv-2508.20296-CC-BY-4.0.txt. Characters: every retained excerpt is pure ASCII, so the delivered engine, XeLaTeX with its default Latin Modern text font, reports no missing character.
\begin{definition} \label{defn:couples}
    
A sequence $(F_n',F_n)$ of pairs of finite subsets
of $G$ with $F_n'\subset F_n$ is called {\it a sequence of F{\o}lner couples}
if there is a constant $C < \infty$ such that for all $n$, we have
 \begin{enumerate}
    \item $\#F_n \leqslant C\#F_n'$
    \item $d_S(F_n',G\setminus F_n) \geqslant n$
\end{enumerate}
\end{definition}

\begin{thm}\label{pr:dimtoisodiam}
Assume that a finitely generated amenable group $G$ has finite asymptotic dimension $d$ with the corresponding control function
 $f(n)$.
Then there exists a sequence of F{\o}lner couples $(F'_n,F_n)$ with $C = d+1$ such that
$\diam(F_n) \leqslant f(2n)+2n$.
\end{thm}

\begin{proof}

We recall that by one of the equivalent definitions of amenability, the group $G$ is amenable if it admits a finitely additive, non-negatively valued,  $G$-invariant measure on all subsets, with total measure one.
One can speak about left-invariant means, which are invariant for the multiplication from the left;
about right-invariant means, which are invariant for the multiplication from the right, and about two-sided invariant means. 
Amenability is equivalent to each of these three conditions: existence of left-invariant means, existence of right-invariant means, or existence of two-sided invariant means.

The following elementary lemma might be reminiscent of the mass transportation argument for amenable groups.
\begin{lemma}\label{lem:lemma1}
Let $G$ be a finitely generated amenable group and let $\mu$ be a right-invariant mean on subsets of $G$. 
Consider a family of disjoint subsets $\{B_{\alpha}\}$ of $G$, $\alpha \in X$. Here the set $X$ is  countable or finite.
Assume that  
the diameters of  $B_{\alpha}$, for all $\alpha \in X$, are uniformly bounded.
Consider a family of subsets $A_{\alpha}$, such that $A_{\alpha} \subseteq B_{\alpha}$ for all $\alpha \in X$. Assume that there exists $\lambda>0$ such that for each $\alpha $ it holds $\#B_{\alpha} \geqslant \lambda \# A_{\alpha}$.
Then  the total mean of the union of $A_{\alpha}$ satisfies:
$$
\mu (\bigsqcup A_{\alpha}) \leqslant 1/\lambda.
$$


\end{lemma}

If the set $X$ is finite, the claim of the lemma is straightforward.  We would have
$$
\mu (\bigsqcup A_{\alpha}) = \sum_{\alpha \in X} \mu(A_\alpha) \le
1/\lambda \sum_{\alpha \in X} \mu(B_\alpha) = 1/\lambda 
\mu (\bigsqcup B_{\alpha}) \le  1/\lambda.
$$


We need to prove the lemma in the case when the $X$ is not necessarily finite, and 
$\mu (\bigsqcup A_{\alpha})$ and  $\mu (\bigsqcup B_{\alpha})$ are not necessarily equal to $\sum_{\alpha \in X} \mu(A_\alpha)$ and respectively
to $\sum_{\alpha \in X} \mu(B_\alpha)$.


\begin{proof}
We consider  family of subsets $A_{\alpha} \subset B_{\alpha} \subset G$ as in the formulation of the lemma.
We say  $\alpha \in X$ and $\beta \in X$ are equivalent if there exists
$g\in G$ such that $A_{\alpha}= gA_{\beta}$ and $B_{\alpha}=g B_{\beta}$.
Since the diameters of the sets $B_{\alpha}$ are uniformly bounded, the number of equivalence classes of our equivalence relation on $X$ is finite.
We enumerate  corresponding sets of indexes 
as $C_1,\dots, C_m$. We have $C_i \subset X$ and $\bigsqcup_{k = 1}^{m}C_k = X$.
We consider total means  of the union of $A_{\alpha}$ and $B_{\alpha}$ in a given class $C_i$:
$$
M^A_i = \mu (\bigsqcup_{ \alpha \in C_i} A_{\alpha}), \hspace{0.5cm}
M^B_i = \mu (\bigsqcup_{\alpha \in C_i} B_{\alpha}).
$$
 
Observe that the cardinality of the sets $A_{\alpha}$ is the same for all
 pairs $(A_{\alpha}, B_{\alpha})$ corresponding to the same equivalence class. We 
denote 
by $a_i$ the cardinality of $A_\alpha$ for $(A_{\alpha}, B_{\alpha})$ corresponding to the class $C_i$. Analogously, the cardinality of the sets
$B_{\alpha}$ for $(A_\alpha, B_{\alpha})$ for the same equivalence class are the same.
We denote  by
$b_i$ the cardinality of $B_{\alpha}$ for $(A_\alpha, B_{\alpha})$
corresponding to the class $C_i$.

 By the assumption of the lemma, we know that $b_i \geqslant \lambda a_i$.
 Fix some $\alpha$ in this class $C_i$.  For each $\beta \in C_i$ fix $g_\beta$ such that $A_{\beta} = g_\beta A_\alpha$ and we assume that $g_\alpha=e$.
For each point $h \in A_\alpha$, we can consider its images by translations
 $g_{\beta} h$, for all $\beta \in C_i$.
We denote this set of points $G_i^h$.
Fix some $h \in A_\alpha$, put $G_i= G_i^h$ and consider $M_i^G = \mu(G_i)$.
Observe that for each $h' \in A_{\alpha}$, we have 
$G_i^{h'} = 
\{g_{\beta} h h^{-1}h'\} = G_i h^{-1}h'$.
Since the right-invariant mean $\mu$ is finitely additive and since $A_\alpha$ is finite, we can 
therefore conclude that  $M_i^A = a_i M_i^G$ and $M_i^B = b_i M_i^G$.
 Hence we know that 
$$
1 \geqslant \mu (\bigsqcup_{\alpha \in X}B_{\alpha}) =
\sum_{i = 1}^{m} M^B_i = \sum_{i = 1}^{m} b_i M^G_i \geqslant
$$
 
$$
\geqslant \lambda \sum_{i = 1}^{m} a_i M^G_i = \lambda \sum_{i = 1}^{m} M^A_i = \lambda \cdot \mu (\bigsqcup_{i = 1}^{\infty} A_i),
$$
 and this completes the proof of the lemma.
\end{proof}


Now we  prove  Theorem \ref{pr:dimtoisodiam}.

Fix $n\ge 1$.
We use the definition of the control function  for the asymptotic dimension for $r=2n$.
We know that  we can color the elements of 
 $G$ in  $d+1$ colors and partition the elements of each color into parts such that the following holds. The distance between distinct parts of the same color is greater than $2n$,
 and each part has diameter at most $f(2n)$.

By the assumption of the theorem,  $G$ is amenable. We fix a right-invariant mean on the subsets of $G$.
Observe that there exists at least one color with the mean of the union of parts
of this color
 $\geqslant \frac{1}{d+1}$. Choose one  such color, and 
  denote the parts of this color by $A_{\alpha}$, $\alpha \in X$, here $X$ is either a countable or finite set.

By our assumption on their diameter, we know that the 
   $n$-neighborhoods  $B(A_{\alpha}, n)$ are pairwise non-intersecting.

We apply Lemma \ref{lem:lemma1} to the parts  $A_{\alpha}$ of the fixed color, defined
above,
and to their neighborhoods  $B_{\alpha}=B(A_{\alpha}, n)$, $n \in \mathbb{N}$.
We conclude that for some $\alpha$, it holds $\#B(A_{\alpha},n) \leqslant (d+1)\#A_\alpha$.


We  put $F'_n = A_{\alpha}$ and $F_n = B(A_\alpha, n)$.
We know that for all $n$
$$
\frac{\# F_n}{\# F'_{n}} \le d+1,
$$
and thus $(F'_n, F_n)$ satisfies the first property in the definition of F{\o}lner couples (see Definition \ref{defn:couples}).

Since, by construction, $F_n$ is a $n$-neighborhood of $F'_n$, it is clear that the distance 
$$
 d (F_n', G \setminus F_n) \ge n,
$$
and thus $(F'_n, F_n)$ also satisfies the second property of Definition \ref{defn:couples}. Therefore  $(F'_n, F_n)$ forms indeed a sequence  of F{\o}lner couples.

Finally, observe that the diameter of $F_n = B(A_\alpha, n)$ is at most $2n$ plus
the diameter of $A_{\alpha}$, and hence the diameter of $F_n$ is at most
$f(2n)+2n$. This completes the proof of the Theorem.
 
\end{proof}

\end{document}
